% Insert position: end of Section 3 (Problem Definition) or beginning of Section 4 (before the PIWM architecture equations).

We adopt the AIDA model~\cite{barry1987development} and the BDI framework~\cite{rao1995bdi} for three reasons. First, AIDA is one of the few sales-funnel schemas that explicitly separates pre-commitment phases (\emph{Attention}, \emph{Interest}, and \emph{Desire}) from the commitment phase (\emph{Action}), making it well-suited to proactive intervention before an explicit purchase decision. Alternative customer-journey schemas such as Howard--Sheth emphasize broader post-purchase satisfaction loops, which are outside our pre-interaction setting. Second, BDI provides a compact tripartite representation that maps naturally to short-horizon visual cues: belief captures inferred uncertainty, desire captures preference signals, and intention captures action readiness. Third, both schemas are textual and schema-constrained, allowing the language model to verbalize intermediate reasoning during training and diagnostic evaluation. Exploring alternative state schemas, including customer-journey models, Howard--Sheth-style decision models, or RFM segmentation, is left for future work; ablations on BDI fields and observable evidence (Sections~6.5.6--6.5.7 after merge) provide initial sensitivity analysis for the chosen schema components.

